\FloatBarrier
\section{本章小结}
\label{sec:condition-perturbation-summary}

本章开头问：只知道误差大小的上限时，由手中的数据算出的结论还有多可靠？答案要先分清误差在哪里，再分清关心的是哪一种输出。

误差在输入向量上、矩阵固定时，线性性给出对任意大小扰动都成立的精确增量$\bm A\bm h$。绝对误差最多放大$\sigma_{\max}(\bm A)$倍，沿最大右奇异方向达到。

相对误差的放大还取决于基点：固定输入处的相对条件数是
\[
 \frac{\lVert\bm A\rVert_2\lVert\bm x\rVert_2}{\lVert\bm A\bm x\rVert_2},
\]
输入落在被压缩的方向上时它大。对输入也取最坏值，得到矩阵条件数
\[
 \kappa_2(\bm A)=\frac{\sigma_{\max}}{\sigma_{\min}},
\]
它衡量最强与最弱作用之差，不衡量矩阵元素的大小，也不描述反复作用的长期趋势。

误差在矩阵本身上时，不同输出有不同的稳定性条件。对称矩阵按序配对的特征值与任意矩阵的奇异值都移动不超过$\lVert\bm E\rVert_2$，不要求谱值彼此分开；据此，区间不含零的特征值保持符号，$\lambda_{\min}(\bm S)>\varepsilon$保证正定，$\sigma_k(\bm A)>\varepsilon$保证秩至少为$k$。

可逆矩阵到奇异矩阵的谱范数距离恰为$\sigma_{\min}$，相对距离为$1/\kappa_2(\bm A)$。精确秩会在任意小的扰动下跳变，可靠的只是给定误差预算下的有效秩。

所选的特征子空间要用主夹角或投影之差比较，单个特征向量的变化不能代表子空间的变化。偏转由扰动与谱隙之比控制：交叉间隔$\gamma>0$时$\sin\theta_{\max}\leq\lVert\bm E\rVert_2/\gamma$，已知原谱隙$\delta$且$\varepsilon<\delta/2$时不超过$2\varepsilon/\delta$。重特征值内部的基不稳定，但只要目标组与其余谱分开，整个子空间稳定；奇异子空间由截断处的奇异值间隙控制。

一般方阵没有对称结构。可对角化时，Bauer--Fike界使扰动后的特征值离原谱不超过$\kappa_2(\bm V)\lVert\bm E\rVert_2$，敏感性由特征基的条件数而不是$\kappa_2(\bm A)$决定；不可对角化时，特征值的移动可以与扰动不成比例，例如$\sqrt t$对$t$。特征方向的偏转由$\lVert(\bm A_{22}-\mu\bm I)^{-1}\rVert_2\lVert\bm E\rVert_2$控制，其中$\bm A_{22}$是所选不变子空间之外的部分；它可以远大于扰动与特征值间隔之比，即使特征值完全不动，特征方向也可能明显偏转。

所以判断一个谱或分解的结果是否可靠，要同时说明对象的结构（对称、可对角化）、所用的分离条件（谱隙、奇异值间隙）和误差预算；仅凭$\kappa_2(\bm A)$不能判断特征值或子空间是否敏感。这些界都是有限变化界，沿指定方向的一阶响应与带余项的近似误差在第~\ref{chap:differential-perturbation-analysis}章讨论。下一章转向同一个矩阵的反复作用：误差每一步都要再经过一次作用，长期结果由谱半径决定。
